\section{Recover the earlier sprite without inventing a new picture}
Expand the existing row bytes from most significant bit to least significant bit. The
first row, \texttt{16}, gives black, black, black, white, black, white, white, black.
Each pixel is opaque. An eight-by-eight output has no per-row padding: the next row
begins thirty-two bytes after the preceding one in RGBA-word memory.

The supplied sixty-one-byte QOI fixture describes that same image. Its fourteen-byte
header gives width eight, height eight, channels four and colorspace zero. The first
payload byte is \texttt{c2}, a run of three initial opaque black pixels. The next,
\texttt{55}, has three DIFF fields of \texttt{01}; each means minus one, turning
zero components into 255 modulo 256. The following \texttt{35} fetches the black cache
entry. These three chunks account for the first five pixels.

Continue by recording output index, input offset, previous pixel and the cache slots
used for black and white. A run may cross a row boundary; the stream counts pixels,
not row records. Divide the output index by eight only when explaining its row and
column. The decoder does not need to insert a new predictor state at each row.

\begin{figure}[htbp]
\centering
\begin{tikzpicture}[x=0.55cm,y=0.55cm,alt={The eight-by-eight monochrome sprite. Rows, from top to bottom, are hexadecimal 16, 3c, 7e, db, ff, 24, 42 and 81. In each row, bit seven is the leftmost square; a set bit is white and a clear bit black. All pixels are opaque.}]
\foreach \row/\value in {0/22,1/60,2/126,3/219,4/255,5/36,6/66,7/129}{
  \foreach \col in {0,...,7}{
    \pgfmathtruncatemacro{\lit}{mod(floor(\value/pow(2,7-\col)),2)}
    \ifnum\lit=1\relax
      \filldraw[fill=white,draw=gray] (\col,-\row) rectangle +(1,-1);
    \else
      \filldraw[fill=black,draw=gray] (\col,-\row) rectangle +(1,-1);
    \fi
  }
}
\end{tikzpicture}
\caption{The logical source sprite from chapter three. The row codes in the description provide a nonvisual specification; this drawing is not an IDE capture.}
\end{figure}
